\documentclass[11pt]{article}
\usepackage[utf8]{inputenc}
\usepackage{amsmath,amssymb,booktabs,geometry,verbatim}
\usepackage{CJKutf8}
\geometry{margin=1in}
\usepackage[hidelinks]{hyperref}

\renewcommand{\labelitemi}{$\bullet$}
\newcommand{\cn}[1]{\begin{CJK}{UTF8}{gbsn}#1\end{CJK}}

\title{CSP-J 2025 (Junior)}
\date{2025-10-25, 8:30--12:00 (3.5 hours)}

\begin{document}
\maketitle

\noindent
In every subtask table, a cell that reads ``as above'' repeats the entry of
the row above it (the original tables merge such cells).

\clearpage
\section*{Problem 1: Forming a Number (\cn{拼数})
  \hfill{\normalsize\href{https://www.luogu.com.cn/problem/P14357}{Luogu P14357}}}

\noindent
Time limit: 1.0 s.
Memory limit: 512 MB.

\medskip
\noindent
Little R is learning string processing.
Little X gives Little R a string $s$,
which consists only of lowercase English letters and digits,
and which \textbf{contains at least one digit among $1 \sim 9$}.
Little X wants Little R to take any number of digits of $s$ and
concatenate them, in any order, into a positive integer.
\textbf{Note: Little R may choose equal digits of $s$,
but every digit (character) of $s$ may be used at most once.}
For example, if $s$ is \texttt{1a01b},
then Little R may choose the 1st, 3rd and 4th characters simultaneously,
which are $1, 0, 1$ respectively, and form the positive integer $101$ or
the positive integer $110$;
but Little R cannot form the positive integer $111$,
because $s$ contains only two digits $1$.
Little R wants to know which of all the positive integers he can form is
the largest.
You need to help Little R find the maximum positive integer he can form.

\subsection*{Input}

The first line of input contains a string $s$,
the string Little X gave to Little R.

\subsection*{Output}

Output one line with one positive integer,
the maximum positive integer that Little R can form.

\subsection*{Sample 1 input}
\begin{verbatim}
5
\end{verbatim}

\subsection*{Sample 1 output}
\begin{verbatim}
5
\end{verbatim}

\subsection*{Sample 2 input}
\begin{verbatim}
290es1q0
\end{verbatim}

\subsection*{Sample 2 output}
\begin{verbatim}
92100
\end{verbatim}

\subsection*{Explanation of sample 2}

The string $s$ contains the digits $2, 9, 0, 1, 0$.
It can be shown that the maximum positive integer Little R can form is
$92100$.

\subsection*{Sample 3}

See the files \texttt{number/number3.in} and \texttt{number/number3.ans}
in the contestant's directory.
This sample satisfies the constraints of test cases $9 \sim 11$.

\subsection*{Sample 4}

See the files \texttt{number/number4.in} and \texttt{number/number4.ans}
in the contestant's directory.
This sample satisfies the constraints of test case $20$.

\subsection*{Constraints}

Let $|s|$ denote the length of the string $s$.
For all test data it is guaranteed that:
\begin{itemize}
\item $1 \leq |s| \leq 10^6$;
\item $s$ consists only of lowercase English letters and digits,
and contains at least one digit among $1 \sim 9$.
\end{itemize}

\begin{center}
\begin{tabular}{ccc}
\toprule
Test case & $|s| \leq$ & Special property \\
\midrule
$1$ & $1$ & A \\
$2$ & $2$ & as above \\
$3$ & as above & none \\
$4$ & $10$ & A \\
$5, 6$ & as above & none \\
$7, 8$ & $10^2$ & A \\
$9 \sim 11$ & as above & none \\
$12$ & $10^3$ & A \\
$13, 14$ & as above & none \\
$15$ & $10^5$ & A \\
$16, 17$ & as above & B \\
$18, 19$ & as above & none \\
$20$ & $10^6$ & A \\
$21, 22$ & as above & B \\
$23 \sim 25$ & as above & none \\
\bottomrule
\end{tabular}
\end{center}

\begin{itemize}
\item Special property A: $s$ consists only of digits.
\item Special property B: $s$ contains at most $10^3$ digits.
\end{itemize}

\clearpage
\section*{Problem 2: Seats (\cn{座位})
  \hfill{\normalsize\href{https://www.luogu.com.cn/problem/P14358}{Luogu P14358}}}

\noindent
Time limit: 1.0 s.
Memory limit: 512 MB.

\medskip
\noindent
The second round of CSP-J 2025 is in progress.
The examination room of Little R has $n \times m$ contestants in total,
and the CSP-J 2025 first-round scores of all contestants are
\textbf{pairwise distinct}.
All $n \times m$ contestants are assigned seats according to their
CSP-J 2025 first-round scores,
\textbf{from high to low in a snake-like order},
arranged in $n$ \textbf{rows} and $m$ \textbf{columns}.
Concretely, let the scores of all contestants in Little R's examination
room, \textbf{from high to low}, be $s_1 > s_2 > \dots > s_{n \times m}$.
Then the contestant with score $s_1$ sits in \textbf{column} $1$,
\textbf{row} $1$;
the contestant with score $s_2$ sits in \textbf{column} $1$,
\textbf{row} $2$;
$\dots$;
the contestant with score $s_n$ sits in \textbf{column} $1$,
\textbf{row} $n$;
the contestant with score $s_{n+1}$ sits in \textbf{column} $2$,
\textbf{row} $n$;
$\dots$;
the contestant with score $s_{2n}$ sits in \textbf{column} $2$,
\textbf{row} $1$;
the contestant with score $s_{2n+1}$ sits in \textbf{column} $3$,
\textbf{row} $1$;
and so on.

For example, if $n = 4, m = 5$,
then all $4 \times 5 = 20$ contestants are assigned seats in decreasing
order of their CSP-J 2025 first-round scores,
following the arrows in the figure below.

\begin{center}
\setlength{\unitlength}{1mm}
\begin{picture}(130,92)(0,0)
\put(29,88){\makebox(0,0){Column 1}}
\put(51,88){\makebox(0,0){Column 2}}
\put(73,88){\makebox(0,0){Column 3}}
\put(95,88){\makebox(0,0){Column 4}}
\put(117,88){\makebox(0,0){Column 5}}
\put(0,73){\makebox(0,0)[l]{Row 1}}
\put(0,51){\makebox(0,0)[l]{Row 2}}
\put(0,29){\makebox(0,0)[l]{Row 3}}
\put(0,7){\makebox(0,0)[l]{Row 4}}
\multiput(22,0)(22,0){5}{\multiput(0,0)(0,22){4}{\framebox(14,14){}}}
\put(29,73){\line(0,-1){66}}
\put(29,7){\line(1,0){22}}
\put(51,7){\line(0,1){66}}
\put(51,73){\line(1,0){22}}
\put(73,73){\line(0,-1){66}}
\put(73,7){\line(1,0){22}}
\put(95,7){\line(0,1){66}}
\put(95,73){\line(1,0){22}}
\put(117,73){\vector(0,-1){64}}
\end{picture}
\end{center}

Given the number of \textbf{rows} $n$ and the number of \textbf{columns}
$m$ of the seats in Little R's examination room,
together with the CSP-J 2025 first-round scores
$a_1, a_2, \dots, a_{n \times m}$ of all contestants in Little R's
examination room, where $a_1$ is the CSP-J 2025 first-round score of
Little R,
you need to help Little R find out in which \textbf{column} and which
\textbf{row} his seat is.

\subsection*{Input}

The first line of input contains two positive integers $n, m$,
the number of \textbf{rows} and the number of \textbf{columns} of the seats
in Little R's examination room.

The second line of input contains $n \times m$ positive integers
$a_1, a_2, \dots, a_{n \times m}$,
the CSP-J 2025 first-round scores of all contestants in Little R's
examination room, where $a_1$ is the CSP-J 2025 first-round score of
Little R.

\subsection*{Output}

Output one line with two positive integers $c, r$,
meaning that Little R's seat is in \textbf{column} $c$, \textbf{row} $r$.

\subsection*{Sample 1 input}
\begin{verbatim}
2 2
99 100 97 98
\end{verbatim}

\subsection*{Sample 1 output}
\begin{verbatim}
1 2
\end{verbatim}

\subsection*{Explanation of sample 1}

In decreasing order of score,
the contestant with score $100$ sits in \textbf{column} $1$,
\textbf{row} $1$;
the contestant with score $99$ sits in \textbf{column} $1$,
\textbf{row} $2$;
the contestant with score $98$ sits in \textbf{column} $2$,
\textbf{row} $2$;
the contestant with score $97$ sits in \textbf{column} $2$,
\textbf{row} $1$.
Little R's score is $99$,
so his seat is in \textbf{column} $1$, \textbf{row} $2$.

\subsection*{Sample 2 input}
\begin{verbatim}
2 2
98 99 100 97
\end{verbatim}

\subsection*{Sample 2 output}
\begin{verbatim}
2 2
\end{verbatim}

\subsection*{Explanation of sample 2}

In decreasing order of score,
the contestant with score $100$ sits in \textbf{column} $1$,
\textbf{row} $1$;
the contestant with score $99$ sits in \textbf{column} $1$,
\textbf{row} $2$;
the contestant with score $98$ sits in \textbf{column} $2$,
\textbf{row} $2$;
the contestant with score $97$ sits in \textbf{column} $2$,
\textbf{row} $1$.
Little R's score is $98$,
so his seat is in \textbf{column} $2$, \textbf{row} $2$.

\subsection*{Sample 3 input}
\begin{verbatim}
3 3
94 95 96 97 98 99 100 93 92
\end{verbatim}

\subsection*{Sample 3 output}
\begin{verbatim}
3 1
\end{verbatim}

\subsection*{Constraints}

For all test data it is guaranteed that:
\begin{itemize}
\item $1 \leq n \leq 10$, $1 \leq m \leq 10$;
\item for all $1 \leq i \leq n \times m$ we have $1 \leq a_i \leq 100$,
and $a_1, a_2, \dots, a_{n \times m}$ are pairwise distinct.
\end{itemize}

\begin{center}
\begin{tabular}{cccc}
\toprule
Test case & $n \leq$ & $m \leq$ & Special property \\
\midrule
$1$ & $1$ & $1$ & AB \\
$2, 3$ & as above & $10$ & none \\
$4, 5$ & $10$ & $1$ & as above \\
$6$ & $2$ & $2$ & A \\
$7$ & as above & as above & B \\
$8, 9$ & as above & as above & none \\
$10$ & as above & $10$ & A \\
$11$ & as above & as above & B \\
$12 \sim 14$ & as above & as above & none \\
$15 \sim 17$ & $10$ & $2$ & as above \\
$18 \sim 20$ & as above & $10$ & as above \\
\bottomrule
\end{tabular}
\end{center}

Special property A: for all $1 \leq i \leq n \times m$ we have $a_i = i$.

Special property B: for all $1 \leq i \leq n \times m$ we have
$a_i = n \times m - i + 1$.

\clearpage
\section*{Problem 3: XOR Sum (\cn{异或和})
  \hfill{\normalsize\href{https://www.luogu.com.cn/problem/P14359}{Luogu P14359}}}

\noindent
Time limit: 1.0 s.
Memory limit: 512 MB.

\medskip
\noindent
Little R has a sequence of $n$ non-negative integers $a_1, a_2, \dots, a_n$.
The weight of an interval $[l, r]$ ($1 \leq l \leq r \leq n$) is defined
as the bitwise XOR of $a_l, a_{l+1}, \dots, a_r$,
that is, $a_l \oplus a_{l+1} \oplus \dots \oplus a_r$,
where $\oplus$ denotes bitwise exclusive or in binary.

Little X gives Little R a non-negative integer $k$.
Little X wants Little R to choose as many \textbf{pairwise disjoint}
intervals of the sequence as possible,
such that the weight of every chosen interval equals $k$.
Two intervals $[l_1, r_1], [l_2, r_2]$ intersect if and only if they
both contain at least one common index,
that is, there exists $1 \leq i \leq n$ with $l_1 \leq i \leq r_1$ and
$l_2 \leq i \leq r_2$.

For example, for the sequence $[2, 1, 0, 3]$,
if $k = 2$, then Little R may choose the interval $[1, 1]$ and the
interval $[2, 4]$, whose weights are $2$ and $1 \oplus 0 \oplus 3 = 2$
respectively;
if $k = 3$, then Little R may choose the interval $[1, 2]$ and the
interval $[4, 4]$, whose weights are $1 \oplus 2 = 3$ and $3$
respectively.

You need to help Little R find the maximum number of intervals he can
choose.

\subsection*{Input}

The first line of input contains two non-negative integers $n, k$,
the length of Little R's sequence and the non-negative integer that
Little X gave to Little R.

The second line of input contains $n$ non-negative integers
$a_1, a_2, \dots, a_n$, Little R's sequence.

\subsection*{Output}

Output one line with one non-negative integer,
the maximum number of intervals that Little R can choose.

\subsection*{Sample 1 input}
\begin{verbatim}
4 2
2 1 0 3
\end{verbatim}

\subsection*{Sample 1 output}
\begin{verbatim}
2
\end{verbatim}

\subsection*{Explanation of sample 1}

Little R may choose the interval $[1, 1]$ and the interval $[2, 4]$,
whose XOR sums are $2$ and $1 \oplus 0 \oplus 3 = 2$ respectively.
It can be shown that the maximum number of intervals Little R can choose
is $2$.

\subsection*{Sample 2 input}
\begin{verbatim}
4 3
2 1 0 3
\end{verbatim}

\subsection*{Sample 2 output}
\begin{verbatim}
2
\end{verbatim}

\subsection*{Explanation of sample 2}

Little R may choose the interval $[1, 2]$ and the interval $[4, 4]$,
whose XOR sums are $1 \oplus 2 = 3$ and $3$ respectively.
It can be shown that the maximum number of intervals Little R can choose
is $2$.

\subsection*{Sample 3 input}
\begin{verbatim}
4 0
2 1 0 3
\end{verbatim}

\subsection*{Sample 3 output}
\begin{verbatim}
1
\end{verbatim}

\subsection*{Explanation of sample 3}

Little R may choose the interval $[3, 3]$, whose XOR sum is $0$.
It can be shown that the maximum number of intervals Little R can choose
is $1$.
Note: Little R cannot choose the interval $[3, 3]$ and the interval
$[1, 4]$ at the same time, because both intervals contain the index $3$.

\subsection*{Sample 4}

See the files \texttt{xor/xor4.in} and \texttt{xor/xor4.ans} in the
contestant's directory.

This sample satisfies the constraints of test cases $4, 5$.

\subsection*{Sample 5}

See the files \texttt{xor/xor5.in} and \texttt{xor/xor5.ans} in the
contestant's directory.

This sample satisfies the constraints of test cases $9, 10$.

\subsection*{Sample 6}

See the files \texttt{xor/xor6.in} and \texttt{xor/xor6.ans} in the
contestant's directory.

This sample satisfies the constraints of test cases $14, 15$.

\subsection*{Constraints}

For all test data it is guaranteed that:
\begin{itemize}
\item $1 \leq n \leq 5 \times 10^5$, $0 \leq k < 2^{20}$;
\item for all $1 \leq i \leq n$ we have $0 \leq a_i < 2^{20}$.
\end{itemize}

\begin{center}
\begin{tabular}{cccc}
\toprule
Test case & $n \leq$ & $k$ & Special property \\
\midrule
$1$ & $2$ & $= 0$ & A \\
$2$ & $10$ & $\leq 1$ & B \\
$3$ & $10^2$ & $= 0$ & A \\
$4, 5$ & as above & $\leq 1$ & B \\
$6 \sim 8$ & as above & $\leq 255$ & C \\
$9, 10$ & $10^3$ & as above & as above \\
$11, 12$ & as above & $< 2^{20}$ & none \\
$13$ & $2 \times 10^5$ & $\leq 1$ & B \\
$14, 15$ & as above & $\leq 255$ & C \\
$16$ & as above & $< 2^{20}$ & none \\
$17$ & $5 \times 10^5$ & $\leq 255$ & C \\
$18 \sim 20$ & as above & $< 2^{20}$ & none \\
\bottomrule
\end{tabular}
\end{center}

Special property A: for all $1 \leq i \leq n$ we have $a_i = 1$.

Special property B: for all $1 \leq i \leq n$ we have $0 \leq a_i \leq 1$.

Special property C: for all $1 \leq i \leq n$ we have $0 \leq a_i \leq 255$.

\clearpage
\section*{Problem 4: Polygon (\cn{多边形})
  \hfill{\normalsize\href{https://www.luogu.com.cn/problem/P14360}{Luogu P14360}}}

\noindent
Time limit: 1.0 s.
Memory limit: 512 MB.

\medskip
\noindent
Little R likes playing with small sticks.
Little R has $n$ sticks, and the $i$-th ($1 \leq i \leq n$) stick has
length $a_i$.

Little X wants Little R to choose some of these $n$ sticks and join them
end to end, in any order, into a polygon.
Little R does not know the condition under which sticks can form a
polygon, so Little X tells him the condition directly:
$m$ sticks of lengths $l_1, l_2, \dots, l_m$ can form a polygon if and
only if $m \geq 3$ and the sum of the lengths of all the sticks is
\textbf{strictly greater than} twice the maximum length of the sticks,
that is, $\sum_{i=1}^{m} l_i > 2 \times \max_{i=1}^{m} l_i$.

Now that Little R knows the condition for sticks to form a polygon,
Little X poses a harder question:
how many ways are there to choose sticks such that the chosen sticks
can form a polygon?
You need to help Little R count the number of ways to choose sticks such
that the chosen sticks can form a polygon.
Two ways are different if and only if the \textbf{sets of indices} of the
chosen sticks are different,
that is, there exists $1 \leq i \leq n$ such that one of the ways chooses
the $i$-th stick but the other does not.
Since the answer may be large, you only need to output the answer modulo
$998,244,353$.

\subsection*{Input}

The first line of input contains a positive integer $n$,
the number of Little R's sticks.

The second line of input contains $n$ positive integers
$a_1, a_2, \dots, a_n$, the lengths of Little R's sticks.

\subsection*{Output}

Output one line with one non-negative integer,
the number of ways to choose sticks such that the chosen sticks can form
a polygon, modulo $998,244,353$.

\subsection*{Sample 1 input}
\begin{verbatim}
5
1 2 3 4 5
\end{verbatim}

\subsection*{Sample 1 output}
\begin{verbatim}
9
\end{verbatim}

\subsection*{Explanation of sample 1}

There are the following $9$ ways to choose sticks such that the chosen
sticks can form a polygon:
\begin{enumerate}
\item choose sticks $2, 3, 4$; the total length is $2 + 3 + 4 = 9$ and
the maximum length is $4$;
\item choose sticks $2, 4, 5$; the total length is $2 + 4 + 5 = 11$ and
the maximum length is $5$;
\item choose sticks $3, 4, 5$; the total length is $3 + 4 + 5 = 12$ and
the maximum length is $5$;
\item choose sticks $1, 2, 3, 4$; the total length is $1 + 2 + 3 + 4 = 10$
and the maximum length is $4$;
\item choose sticks $1, 2, 3, 5$; the total length is $1 + 2 + 3 + 5 = 11$
and the maximum length is $5$;
\item choose sticks $1, 2, 4, 5$; the total length is $1 + 2 + 4 + 5 = 12$
and the maximum length is $5$;
\item choose sticks $1, 3, 4, 5$; the total length is $1 + 3 + 4 + 5 = 13$
and the maximum length is $5$;
\item choose sticks $2, 3, 4, 5$; the total length is $2 + 3 + 4 + 5 = 14$
and the maximum length is $5$;
\item choose sticks $1, 2, 3, 4, 5$; the total length is
$1 + 2 + 3 + 4 + 5 = 15$ and the maximum length is $5$.
\end{enumerate}

\subsection*{Sample 2 input}
\begin{verbatim}
5
2 2 3 8 10
\end{verbatim}

\subsection*{Sample 2 output}
\begin{verbatim}
6
\end{verbatim}

\subsection*{Explanation of sample 2}

There are the following $6$ ways to choose sticks such that the chosen
sticks can form a polygon:
\begin{enumerate}
\item choose sticks $1, 2, 3$; the total length is $2 + 2 + 3 = 7$ and
the maximum length is $3$;
\item choose sticks $3, 4, 5$; the total length is $3 + 8 + 10 = 21$ and
the maximum length is $10$;
\item choose sticks $1, 2, 4, 5$; the total length is $2 + 2 + 8 + 10 = 22$
and the maximum length is $10$;
\item choose sticks $1, 3, 4, 5$; the total length is $2 + 3 + 8 + 10 = 23$
and the maximum length is $10$;
\item choose sticks $2, 3, 4, 5$; the total length is $2 + 3 + 8 + 10 = 23$
and the maximum length is $10$;
\item choose sticks $1, 2, 3, 4, 5$; the total length is
$2 + 2 + 3 + 8 + 10 = 25$ and the maximum length is $10$.
\end{enumerate}

\subsection*{Sample 3}

See the files \texttt{polygon/polygon3.in} and
\texttt{polygon/polygon3.ans} in the contestant's directory.

This sample satisfies the constraints of test cases $7 \sim 10$.

\subsection*{Sample 4}

See the files \texttt{polygon/polygon4.in} and
\texttt{polygon/polygon4.ans} in the contestant's directory.

This sample satisfies the constraints of test cases $11 \sim 14$.

\subsection*{Constraints}

For all test data it is guaranteed that:
\begin{itemize}
\item $3 \leq n \leq 5\,000$;
\item for all $1 \leq i \leq n$ we have $1 \leq a_i \leq 5\,000$.
\end{itemize}

\begin{center}
\begin{tabular}{ccc}
\toprule
Test case & $n \leq$ & $\max_{i=1}^{n} a_i \leq$ \\
\midrule
$1 \sim 3$ & $3$ & $10$ \\
$4 \sim 6$ & $10$ & $10^2$ \\
$7 \sim 10$ & $20$ & as above \\
$11 \sim 14$ & $500$ & as above \\
$15 \sim 17$ & as above & $1$ \\
$18 \sim 20$ & $5\,000$ & as above \\
$21 \sim 25$ & as above & $5\,000$ \\
\bottomrule
\end{tabular}
\end{center}

\section*{Source}

The statements were translated from the official problems as published
on Luogu (CSP-J 2025, second round, held on 2025-10-25):
\begin{itemize}
\item Problem 1: \texttt{https://www.luogu.com.cn/problem/P14357}
\item Problem 2: \texttt{https://www.luogu.com.cn/problem/P14358}
\item Problem 3: \texttt{https://www.luogu.com.cn/problem/P14359}
\item Problem 4: \texttt{https://www.luogu.com.cn/problem/P14360}
\end{itemize}

\end{document}
